%%%%%%%%%%%%%%%%%%%%%%%%%%%%%%%%%%%%%%%%%%%%%%%%%%%%%%%
% DecL Cheat Sheet  (multi-page)
%
% Created by Stephen J Mildenhall
% (c) 2023
%
% Updated 2026 for aggregate v1.0.0a169 (REFACTOR branch):
%   * grammar source of truth is aggregate/decl.lark (Lark Earley);
%     every form on this card was parse-checked against the live parser
%   * grown 3 -> 4 pages:
%       p1  agg   (the compound-distribution language, incl. renewal wait)
%       p2  reinsurance economics + pnl / xpnl
%       p3  port + the bivariate family (copula, dbvsev, view pairs, clash)
%       p4  distortion
%   * corrections to the a66 card: xps goes with (d|c)histogram, not a
%     parametric severity; a reflected severity needs ssev, not sev;
%     wtdtvar has NO DecL number-list form; variable rating is
%     aggregate-reinsurance only in this release; `*` is severity scaling,
%     not general arithmetic
%
%%%%%%%%%%%%%%%%%%%%%%%%%%%%%%%%%%%%%%%%%%%%%%%%%%%%%%%

\input{cheat_sheet_macros.tex}
\title{DecL Cheat Sheet}

% color scheme defeined here - just change the suffixes
\colorlet{highlightcolor}{highlightcolora}
\colorlet{washedcolor}{washedcolora}
\colorlet{textcolor}{texta}


\begin{document}

%%%%%%%%%%%%%%%%%%%%%%%%%%%%%%%%%%%%%%%%%%%%%%%%%%%%%%%%%%%%%%%%%%%%%%%
% PAGE 1 : agg
%%%%%%%%%%%%%%%%%%%%%%%%%%%%%%%%%%%%%%%%%%%%%%%%%%%%%%%%%%%%%%%%%%%%%%%

{\huge{\bf DecL Cheat Sheet \textnormal{(1/4)}: \texttt{agg}}}

\raggedright An \texttt{Aggregate} compound distribution has the clause structure: \\
\texttt{agg <NAME> <LABEL*> <EXPOSURE> <LAYERS*> <SEVERITY> <OCC\_RE*> <FREQUENCY> <AGG\_RE*> <APPROX*> <ORIENTATION*> <TRAILER*>} \\
Key: \texttt{<INPUT>} user input; \texttt{lower\_case} a DecL keyword; \texttt{CLAUSE} a sub-clause; options \texttt{a|b|c}; \texttt{inf} infinity; \texttt{*} optional. Build with \texttt{build('agg ...')}.

\begin{multicols*}{3}

%--- Name ---
\begin{tikzpicture}
\node [mybox] (box){\begin{minipage}{0.3\textwidth}\raggedright
\texttt{agg <NAME> ...} \\
\texttt{sev <NAME> <SEVERITY>} {\it standalone severity} \\
\texttt{... as <LABEL>}, \texttt{... as "<LABEL>"} {\it display label}

\medskip\it Names match \rm\texttt{[a-zA-Z][.\_:\textasciitilde a-zA-Z0-9-]*}\it{} and are the identity handle; \texttt{as} adds presentation text. A label may also attach to the exposure, a layer, a severity, a reinsurance layer, and an expense group.

\medskip\it Reference: \rm\texttt{agg.<NAME>}, \texttt{sev.<NAME>}, \texttt{port.<NAME>}, \texttt{dist.<NAME>}. \it Clone / transform: \rm\texttt{agg <NEW> agg.<OLD>}, \texttt{2 * agg.X} {\it (homog.)}, \texttt{2 @ agg.X} {\it (inhomog.)}, \texttt{agg.X + 100}.
\end{minipage}};
\node[fancytitle] at (box.north west) {1. Name, labels \& references};
\end{tikzpicture}

%--- Exposure ---
\begin{tikzpicture}
\node [mybox] (box){\begin{minipage}{0.3\textwidth}\raggedright
\texttt{<EL> loss} \\
\texttt{<PREMIUM> premium at <LR> lr} \\
\texttt{<EXPOSURE> exposure at <RATE> rate} \\
\texttt{<CLAIMS> claims} \\
\texttt{dfreq <OUTCOMES> <PROBS*>} {\it fixed empirical freq} \\
\texttt{<T> years <at <RATE> rate>*} {\it renewal horizon}

\medskip\it Outcomes \texttt{[1 2 3 4]} or \texttt{[2:10:2]}; probs \texttt{[.5 .25 .25]} or omitted $=$ equally likely. \texttt{years} pairs only with \texttt{wait} (box 7); its \texttt{rate} is informational and never sizes the count.
\end{minipage}};
\node[fancytitle] at (box.north west) {2. Exposure};
\end{tikzpicture}

%--- Limit ---
\begin{tikzpicture}
\node [mybox] (box){\begin{minipage}{0.3\textwidth}\raggedright
\texttt{<LIMIT> xs <ATTACHMENT> <LABEL*>} \\
\texttt{tower [<BREAKS>]}

\medskip\it Occurrence limits on ground-up severity, unlimited reinstatements, conditional on attaching by default.
\end{minipage}};
\node[fancytitle] at (box.north west) {3. Limit / layers (optional)};
\end{tikzpicture}

\columnbreak

%--- Severity continuous ---
\begin{tikzpicture}
\node [mybox] (box){\begin{minipage}{0.3\textwidth}\raggedright
\texttt{sev <DIST> <MEAN> cv <CV>} \\
\texttt{sev <DIST> <SHAPE1> <SHAPE2*>} \\
\texttt{sev sev.<NAME>} {\it built-in}

\medskip\texttt{<SCALE> * SEV + <LOC>} {\it scale \& shift} \\
\texttt{ssev <NUM> - SEV} {\it reflect (needs} \texttt{ssev}) \\
\texttt{SEV wts [<W>]}, \texttt{SEV wts=<N>} {\it mixture} \\
\texttt{SEV splice [<LB> <UB>]} {\it conditional in layer} \\
\texttt{SEV !} {\it unconditional (when} \texttt{ATTACH}>0) \\
\texttt{ssev ...} {\it spliced / signed severity}
\end{minipage}};
\node[fancytitle] at (box.north west) {4. Severity (continuous)};
\end{tikzpicture}

%--- Severity discrete ---
\begin{tikzpicture}
\node [mybox] (box){\begin{minipage}{0.3\textwidth}\raggedright
\texttt{dsev <OUTCOMES> <PROBS*>}, \texttt{dsev ... !} \\
\texttt{sev (d|c)histogram xps [<X>] [<P>]} \\
\texttt{picks [<X>] [<P>]} {\it match layer picks}

\medskip\it \texttt{dsev [1:6]} is a fair die; \texttt{xps} gives outcomes and probabilities directly, and attaches only to a histogram.
\end{minipage}};
\node[fancytitle] at (box.north west) {5. Severity (discrete)};
\end{tikzpicture}

%--- Frequency ---
\begin{tikzpicture}
\node [mybox] (box){\begin{minipage}{0.3\textwidth}\raggedright
\texttt{poisson}, \texttt{bernoulli}, \texttt{fixed}, \texttt{geometric}, \texttt{logarithmic}, \texttt{binomial <P>}, \texttt{negbin <VAR\_MULT>}, \texttt{neymana <CLAIMS/OCC>}, \texttt{pascal <CV> <CLAIMS/OCC>}

\medskip\texttt{mixed <MIX> <SHAPE1> <SHAPE2*>} \\
\texttt{MIX = gamma|delaporte|ig|sig|sichel|beta}

\medskip\texttt{FREQ zt} {\it zero-truncated};
\texttt{FREQ zm <P0>} {\it zero-modified,} $\Pr(N{=}0){=}p_0$;
{\it trailing} \texttt{!} {\it pins the realized mean to the exposure clause instead of the base mean}
\end{minipage}};
\node[fancytitle] at (box.north west) {6. Frequency};
\end{tikzpicture}

%--- Renewal / wait ---
\begin{tikzpicture}
\node [mybox] (box){\begin{minipage}{0.3\textwidth}\raggedright
\texttt{wait <SEVERITY> <LABEL*>} \\
\texttt{wait <Y> xs <A> <SEVERITY> <LABEL*>} \\
\texttt{dwait <OUTCOMES> <PROBS*>}, \texttt{dwait ... !}

\medskip\it The Sparre-Andersen alternative to a frequency clause: the count comes from the waiting-time law over the \texttt{years} horizon. \texttt{wait} takes the full severity mini-language on the time axis. The layered form is conditional by default; \texttt{!} collapses $\Pr(W \le a)$ into a zero-wait cluster atom.

\medskip\rm\texttt{agg R 10 years sev lognorm 50 cv 2 wait expon 1}
\end{minipage}};
\node[fancytitle] at (box.north west) {7. Renewal waiting time};
\end{tikzpicture}

\columnbreak

%--- Reinsurance ---
\begin{tikzpicture}
\node [mybox] (box){\begin{minipage}{0.3\textwidth}\raggedright
\texttt{occurrence (net of|ceded to) LAYERS} \\
\texttt{aggregate (net of|ceded to) LAYERS}

\medskip\texttt{LAYER = <LIMIT> xs <ATTACH>} \\
\hspace*{1em}\texttt{| <SHARE> so <LIMIT> xs <ATTACH>} \\
\hspace*{1em}\texttt{| <PCT> po <LIMIT> xs <ATTACH>} \\
\hspace*{1em}\texttt{| <FRAC> of <LIMIT> xs <ATTACH>} \\
\texttt{LAYERS = LAYER and LAYER ...} {\it or} \texttt{tower [...]}

\medskip\it \texttt{so} $=$ share of, \texttt{po} $=$ part of, \texttt{of} $=$ fraction of; $0\le\mathrm{share}\le1$. Occurrence applies per claim, aggregate to the annual total. Layers may also carry economics and a label; see page 2.
\end{minipage}};
\node[fancytitle] at (box.north west) {8. Reinsurance};
\end{tikzpicture}

%--- Approximate + orientation ---
\begin{tikzpicture}
\node [mybox] (box){\begin{minipage}{0.3\textwidth}\raggedright
\texttt{approximate exact} {\it (default)}, \texttt{approximate sgamma} {\it shifted gamma}, \texttt{approximate slognorm} {\it shifted lognormal}

\medskip\it Alias \texttt{approx}. Replaces the freq\,$\times$\,sev convolution with one severity matched to the first three aggregate moments.

\medskip\texttt{loss} {\it (default), } \texttt{payoff}. {\it The orientation suffix, last before the trailer: \texttt{payoff} marks an asset-return primitive (more is better) and flips pricing to the dual distortion.}

\medskip\texttt{tweedie <MEAN> <P> <DISP>} {\it a declaration on its own}
\end{minipage}};
\node[fancytitle] at (box.north west) {9. Approximate \& orientation};
\end{tikzpicture}

%--- Trailer + vectorization + math ---
\begin{tikzpicture}
\node [mybox] (box){\begin{minipage}{0.3\textwidth}\raggedright
{\bf Trailer} \texttt{note\{free text\}}, \texttt{tags\{cat, property\}}, \texttt{hints\{bs=100; log2=17\}}, \texttt{doc\{\{\{markdown\}\}\}}. {\it Order-free, at most one of each; only \texttt{doc} may carry blank lines and fenced code.}

\medskip{\bf Vectorize} \rm\texttt{[1 2 3] claims}; \texttt{[250 250 500] xs [0 250 500]}; \texttt{lognorm [1 3] cv [.7 1.5] wts [.6 .4]}. {\it Vectors broadcast; layers / exposure zip.}

\medskip{\bf Math} \rm\texttt{1/2}, \texttt{3**4}, \texttt{3\textasciicircum4}, \texttt{exp(2)}, \texttt{12.4\%}, \texttt{10\_000}. {\it Minus binds to the number, so} $-4^2=16$; \texttt{*} {\it is severity scaling, not general arithmetic.} {\bf Comments} \rm\texttt{\#} {\it to end of line; statements separate on a blank line or} \texttt{;}.
\end{minipage}};
\node[fancytitle] at (box.north west) {10. Trailer, vectors, math};
\end{tikzpicture}

\makefooter
\end{multicols*}

\clearpage

%%%%%%%%%%%%%%%%%%%%%%%%%%%%%%%%%%%%%%%%%%%%%%%%%%%%%%%%%%%%%%%%%%%%%%%
% PAGE 2 : reinsurance economics + pnl / xpnl
%%%%%%%%%%%%%%%%%%%%%%%%%%%%%%%%%%%%%%%%%%%%%%%%%%%%%%%%%%%%%%%%%%%%%%%

{\huge{\bf DecL Cheat Sheet \textnormal{(2/4)}: reinsurance economics, \texttt{pnl} and \texttt{xpnl}}}

\raggedright Page 1 declared the reinsurance \emph{loss} structure. Everything here is \emph{money}: ceded premium, commission, reinstatements, variable rating, and the profit-and-loss statement that books them. \\
A priced layer is \texttt{LAYER <PREMIUM*> <CEDE*> <REINSTATEMENTS*> <VARIABLE*> <LABEL*>}, in that order.

\begin{multicols*}{3}

%--- Ceded premium ---
\begin{tikzpicture}
\node [mybox] (box){\begin{minipage}{0.3\textwidth}\raggedright
\texttt{LAYER deposit <AMOUNT>} \\
\texttt{LAYER rol <RATE>} {\it rate on line} $\times$ share $\times$ limit \\
\texttt{LAYER rate <FRACTION>} {\it of gross premium}

\medskip\texttt{... cede <FRACTION>} {\it ceding commission}

\medskip\it The three premium forms are mutually exclusive. \texttt{cede} is valid only with a premium; the premium is entered \rm gross\it{} of commission.

\medskip\rm\texttt{occurrence net of 250 xs 250 deposit 60 cede 20\%}
\end{minipage}};
\node[fancytitle] at (box.north west) {1. Ceded premium \& commission};
\end{tikzpicture}

%--- Reinstatements ---
\begin{tikzpicture}
\node [mybox] (box){\begin{minipage}{0.3\textwidth}\raggedright
\texttt{reinstatements [0 .5 1 1]} {\it explicit multipliers} \\
\texttt{reinstatements 1 free and 2 at 100\%} \\
\texttt{no reinstatements} {\it one annual limit}

\medskip\it Multipliers $\alpha_j$ are relative to the base rate on line $r$, so a base premium clause is required. \texttt{free} $=$ \texttt{at 0\%}. Counts may be digits or the words \texttt{one}\dots\texttt{five}. $m$ reinstatements $=$ $m+1$ annual limits, so \texttt{no reinstatements} is one limit --- distinct from \rm omitting\it{} the clause, which keeps free and unlimited.

\medskip\it One occurrence layer, one clause. This is the stochastic-ceded property-cat feature.
\end{minipage}};
\node[fancytitle] at (box.north west) {2. Reinstatements};
\end{tikzpicture}

%--- Variable rating ---
\begin{tikzpicture}
\node [mybox] (box){\begin{minipage}{0.3\textwidth}\raggedright
\texttt{swing basic <B> lcm <M> <min <LO>>* <max <HI>>*} \\
\texttt{slide <C1> at <LR1> and <C2> at <LR2> ...} \\
\texttt{pc <SHARE> after <ALLOWANCE>} \\
\texttt{corridor <SHARE> po <WIDTH> xs <ATTACH>}

\medskip\it \texttt{swing} \rm replaces\it{} the deposit / rol / rate premium with the collared affine $\mathrm{clip}(b + m \cdot L_{\text{ceded}}, lo, hi)$; \texttt{slide} \rm replaces\it{} \texttt{cede} with a commission sliding on the ceded loss ratio; \texttt{pc} is a profit commission on the ceded premium; \texttt{corridor} is a loss-ratio band the cedant retains.

\medskip\it At most \rm one\it{} variable feature per program (no stacking), and \rm aggregate reinsurance only\it{} in this release.
\end{minipage}};
\node[fancytitle] at (box.north west) {3. Variable rating};
\end{tikzpicture}

\columnbreak

%--- pnl declaration ---
\begin{tikzpicture}
\node [mybox] (box){\begin{minipage}{0.3\textwidth}\raggedright
\texttt{pnl <NAME> <LABEL*> <PREMIUM> less <ENGINE> <less <EXPENSES>>* <TRAILER*>} \\
\texttt{xpnl <NAME> ...} {\it same syntax}

\medskip\it Profit and loss $=$ premium minus what comes out of the sausage maker. \texttt{less} is a dedicated keyword, never \texttt{-}, so the split cannot collide with severity arithmetic (\texttt{ssev 20 - lognorm}); a \rm second\it{} \texttt{less} is the hard anchor introducing expenses.

\medskip\rm\texttt{pnl P 1000 premium less agg E 700 loss sev lognorm 50 cv 2 poisson}
\end{minipage}};
\node[fancytitle] at (box.north west) {4. \texttt{pnl} / \texttt{xpnl} declaration};
\end{tikzpicture}

%--- Premium head ---
\begin{tikzpicture}
\node [mybox] (box){\begin{minipage}{0.3\textwidth}\raggedright
\texttt{<AMOUNT> premium <LABEL*>} \\
\texttt{inherit premium <LABEL*>} \\
\texttt{derive premium <LABEL*>} \\
\texttt{retro <COLLAR> premium <LABEL*>}

\medskip\it The premium is the \rm consideration\it{} (what you book); the engine's own \texttt{premium at lr} is a separate \rm sizing\it{} input. \texttt{inherit} copies the engine's technical premium (a build error if it has none). \texttt{derive} grosses it up for the expense clause, $P = (T + \text{fixed}) / (1 - \text{premium ratios})$, so premium net of expenses returns exactly $T$; fixed and premium bases only. \texttt{retro} is the account-level retrospective-rating clause, taking the same \texttt{basic ... lcm ... min ... max} collar as \texttt{swing}.

\medskip\rm\texttt{pnl P inherit premium as "gross written" less agg.Book}
\end{minipage}};
\node[fancytitle] at (box.north west) {5. The premium head};
\end{tikzpicture}

%--- Engine source ---
\begin{tikzpicture}
\node [mybox] (box){\begin{minipage}{0.3\textwidth}\raggedright
\texttt{agg <NAME> <AGG-BODY>} {\it inline} \\
\texttt{agg.<NAME>} {\it stored aggregate} \\
\texttt{port.<NAME>} {\it stored portfolio (net-net total)}

\medskip\it A P\&L wraps a \rm complete\it{} stochastic engine, never an anonymous half-body: the inline form reuses the identical \texttt{agg} body from page 1, minus the trailer, which the wrapping \texttt{pnl} owns. So every clause on page 1 is available inside a \texttt{pnl}.
\end{minipage}};
\node[fancytitle] at (box.north west) {6. The engine source};
\end{tikzpicture}

\columnbreak

%--- Expenses ---
\begin{tikzpicture}
\node [mybox] (box){\begin{minipage}{0.3\textwidth}\raggedright
\texttt{less <TERM> and <TERM> ... <LABEL*>} {\it one group} \\
\texttt{less <GROUP> <GROUP> ...} {\it several groups}

\medskip\texttt{TERM = <F> premium expense} \\
\hspace*{1em}\texttt{| <F> loss expense} \\
\hspace*{1em}\texttt{| <AMOUNT> fixed expense}

\medskip\it \texttt{expense} and \texttt{expenses} are both accepted. Terms joined by \texttt{and} \rm combine\it{} into one reported obligation leg; juxtaposed groups stay \rm separate\it{} items. Each group may carry its own \texttt{as} label; the default name is the basis.

\medskip\rm\texttt{less 25\% premium expense as acquisition 1000 fixed expense as overhead}
\end{minipage}};
\node[fancytitle] at (box.north west) {7. Expenses};
\end{tikzpicture}

%--- pnl vs xpnl ---
\begin{tikzpicture}
\node [mybox] (box){\begin{minipage}{0.3\textwidth}\raggedright
\texttt{pnl} {\it the \rm consolidated\it{} single-group net view: one \texttt{Consideration} / \texttt{Obligation} / \texttt{Margin} card.}

\medskip\texttt{xpnl} {\it the \rm exploded\it{} sibling, same syntax: a multi-group walk with one group per step, so you see gross, then each cession, then expenses, as separate blocks that foot to the same margin.}

\medskip\it Both return a \texttt{PnL}. Reach for \texttt{xpnl} to see \rm where\it{} the margin went, for \texttt{pnl} to see \rm what it is\it. Both may be portfolio units.
\end{minipage}};
\node[fancytitle] at (box.north west) {8. \texttt{pnl} vs.\ \texttt{xpnl}};
\end{tikzpicture}

%--- Worked example ---
\begin{tikzpicture}
\node [mybox] (box){\begin{minipage}{0.3\textwidth}\raggedright
\rm\texttt{xpnl Book 10000 premium} \\
\hspace*{1em}\texttt{less agg Book\_e 10000 premium at 75\% lr} \\
\hspace*{2em}\texttt{10000 xs 0} \\
\hspace*{2em}\texttt{sev 10.808 * lognorm 1.75} \\
\hspace*{2em}\texttt{mixed gamma 0.25} \\
\hspace*{2em}\texttt{aggregate net of 5000 xs 10000} \\
\hspace*{3em}\texttt{deposit 1500 pc 50\% after 10\%} \\
\hspace*{1em}\texttt{less 20\% premium expense} \\
\hspace*{1em}\texttt{note\{profit commission on the agg cover\}}

\medskip\it Premium head, engine with an aggregate cession carrying a deposit and a profit commission, then expenses.
\end{minipage}};
\node[fancytitle] at (box.north west) {9. Worked example};
\end{tikzpicture}

\makefooter
\end{multicols*}

\clearpage

%%%%%%%%%%%%%%%%%%%%%%%%%%%%%%%%%%%%%%%%%%%%%%%%%%%%%%%%%%%%%%%%%%%%%%%
% PAGE 3 : port + the bivariate family
%%%%%%%%%%%%%%%%%%%%%%%%%%%%%%%%%%%%%%%%%%%%%%%%%%%%%%%%%%%%%%%%%%%%%%%

{\huge{\bf DecL Cheat Sheet \textnormal{(3/4)}: \texttt{port} and the bivariate family}}

\raggedright These declarations build collections and couplings of aggregates. Each component is an ordinary \texttt{agg} (or \texttt{pnl}) declaration, so everything on pages 1 and 2 is available inside them.

\begin{multicols*}{3}

%--- Portfolio ---
\begin{tikzpicture}
\node [mybox] (box){\begin{minipage}{0.3\textwidth}\raggedright
\texttt{port <NAME> <LABEL*> <TRAILER*>} \\
\hspace*{1em}\texttt{agg <UNIT1> ...} \\
\hspace*{1em}\texttt{agg <UNIT2> ...} \\
\hspace*{1em}\dots

\medskip\it A \texttt{Portfolio} is a list of \rm independent\it{} units sharing a discretization grid; the total is their convolution. Units may be \texttt{agg} or \texttt{pnl} lines, or references \texttt{agg.<NAME>}. Per-unit reinsurance is set on each unit line.

\medskip\it The portfolio trailer sits \rm before\it{} the unit list, so a trailing trailer always binds to the last unit.
\end{minipage}};
\node[fancytitle] at (box.north west) {1. Portfolio (\texttt{port})};
\end{tikzpicture}

%--- Portfolio example ---
\begin{tikzpicture}
\node [mybox] (box){\begin{minipage}{0.3\textwidth}\raggedright
\rm\texttt{port MyBook note\{two lines\}} \\
\hspace*{1em}\texttt{agg A 100 claims sev lognorm 50 cv 2 poisson} \\
\hspace*{1em}\texttt{agg B 200 claims sev gamma 30 cv 1 poisson}

\medskip\it Access units by name (\texttt{p.A}); \texttt{unit\_names\_ex} adds \texttt{total}.
\end{minipage}};
\node[fancytitle] at (box.north west) {2. \texttt{port} example};
\end{tikzpicture}

%--- Bivariate ---
\begin{tikzpicture}
\node [mybox] (box){\begin{minipage}{0.3\textwidth}\raggedright
\texttt{bivariate <NAME> <SHARED-EXPOSURE>} \\
\hspace*{1em}\texttt{agg <A> ...} \\
\hspace*{1em}\texttt{agg <B> ...} \\
\hspace*{1em}\texttt{copula <KIND> <P*>} \\
\hspace*{1em}\texttt{<FREQUENCY*>}

\medskip\it Abbreviation \texttt{bv}. Couples \rm exactly two\it{} component severities by a copula and accumulates them with a \rm shared\it{} outer frequency via a 2D FFT. The shared exposure may be \texttt{<N> claims} or \texttt{dfreq [...]}; the trailing frequency is optional (default poisson). Per-event triggers come from the inner \texttt{dfreq [0 1] [p0 p1]} Bernoulli forms.
\end{minipage}};
\node[fancytitle] at (box.north west) {3. Bivariate (\texttt{bv})};
\end{tikzpicture}

\columnbreak

%--- Copula clause ---
\begin{tikzpicture}
\node [mybox] (box){\begin{minipage}{0.3\textwidth}\raggedright
\texttt{copula <KIND> <P>} {\it one-parameter} \\
\texttt{copula <KIND>} {\it parameter-free} \\
{\it omitted} $=$ \texttt{copula independent}

\medskip\texttt{normal <rho>} $\rho \in (-1, 1)$ \\
\texttt{gumbel <tau>} $\tau \in [0, 1)$ {\it upper tail dependence} \\
\texttt{clayton <tau>} $\tau \in [0, 1)$ {\it lower tail dependence} \\
\texttt{fgm <rho\_s>} $\rho_s \in [-1/3, 1/3]$ \\
\texttt{independent}

\medskip\rm\texttt{bv M 100 claims agg A ... agg B ... copula gumbel 0.4}
\end{minipage}};
\node[fancytitle] at (box.north west) {4. Copula clause};
\end{tikzpicture}

%--- dbvsev ---
\begin{tikzpicture}
\node [mybox] (box){\begin{minipage}{0.3\textwidth}\raggedright
\texttt{dbvsev [<XS>] [<YS>] [[<ROW>] [<ROW>] ...]} \\
\texttt{dbvsev [<XS>] [<YS>]} {\it uniform on the lattice} \\
\texttt{dbvsev [[<X> <Y> <P>] ...]} {\it sparse triples}

\medskip\it The 2-D analogue of \texttt{dsev}: the joint per-claim probability matrix given directly, \texttt{matrix[i][j]} $= \Pr(X{=}x_i, Y{=}y_j)$. Replaces the two \texttt{agg} bodies \rm and\it{} the copula clause. The three surface forms are auto-detected.

\medskip\rm\texttt{bv D 10 claims dbvsev [1 2] [3 4] [[.25 .25] [.25 .25]]}
\end{minipage}};
\node[fancytitle] at (box.north west) {5. Discrete bivariate (\texttt{dbvsev})};
\end{tikzpicture}

\columnbreak

%--- View pairs ---
\begin{tikzpicture}
\node [mybox] (box){\begin{minipage}{0.3\textwidth}\raggedright
\texttt{netceded <AGG>} \\
\texttt{grossceded <AGG>} \\
\texttt{grossnet <AGG>}

\medskip\it One ordinary \texttt{agg} carrying occurrence reinsurance; the keyword names the pair \rm x-then-y\it{} out of \{gross, ceded, net\} (gross leads when present). The two coupled marginals are positions of the \rm same\it{} loss, so their dependence is exact, comonotone within the layer. Lets you see the pair \rm jointly\it, not just their marginals.

\medskip\rm\texttt{netceded agg G 100 claims sev lognorm 50 cv 2 occurrence net of 500 xs 500 poisson}
\end{minipage}};
\node[fancytitle] at (box.north west) {6. View pairs};
\end{tikzpicture}

%--- Clash ---
\begin{tikzpicture}
\node [mybox] (box){\begin{minipage}{0.3\textwidth}\raggedright
\texttt{clash <NAME> <NA> <NB> <NC> claims} \\
\hspace*{1em}\texttt{<LAYERS*> <SEVERITY>} {\it component A} \\
\hspace*{1em}\texttt{<LAYERS*> <SEVERITY>} {\it component B} \\
\hspace*{1em}\texttt{<FREQUENCY*>}

\medskip\it Dependence from interpretable claim counts rather than a copula parameter: $n_a$ claims hit A, $n_b$ hit B, and $n_c$ hit \rm both\it{} (the clash). The solver derives the shared event count $n_0$ and the two per-event Bernoulli triggers $p_a$, $p_b$.

\medskip\rm\texttt{clash C 10 8 2 claims 100 xs 0 sev lognorm 50 cv 2 100 xs 0 sev gamma 40 cv 1 poisson}
\end{minipage}};
\node[fancytitle] at (box.north west) {7. Clash};
\end{tikzpicture}

\makefooter
\end{multicols*}

\clearpage

%%%%%%%%%%%%%%%%%%%%%%%%%%%%%%%%%%%%%%%%%%%%%%%%%%%%%%%%%%%%%%%%%%%%%%%
% PAGE 4 : distortion
%%%%%%%%%%%%%%%%%%%%%%%%%%%%%%%%%%%%%%%%%%%%%%%%%%%%%%%%%%%%%%%%%%%%%%%

{\huge{\bf DecL Cheat Sheet \textnormal{(4/4)}: \texttt{distortion}}}

\raggedright A distortion $g$ defines a coherent risk measure / pricing functional. Declare one by kind and its natural parameters; the kind\,$\to$\,parameter mapping lives on each \texttt{Distortion} subclass (\texttt{decl\_params}), so the grammar holds no per-kind knowledge.

\begin{multicols*}{3}

%--- Distortion decl ---
\begin{tikzpicture}
\node [mybox] (box){\begin{minipage}{0.3\textwidth}\raggedright
\texttt{distortion <NAME> <KIND> <PARAMS> <TRAILER*>}

\medskip\it Alias \texttt{dist}. \texttt{<PARAMS>} is a flat list of the kind's natural parameters, in the order of the table. Reference a stored distortion as \texttt{dist.<NAME>}. A \texttt{dist} line carries the same trailer as every other statement, so it can be described, tagged, and documented.

\medskip\rm\texttt{distortion D1 ph 0.9} \\
\texttt{dist D2 wang 0.25 note\{market standard\}} \\
\texttt{dist D3 power 0 100 2}
\end{minipage}};
\node[fancytitle] at (box.north west) {1. Declaration};
\end{tikzpicture}

%--- Kinds ---
\begin{tikzpicture}
\node [mybox] (box){\begin{minipage}{0.3\textwidth}\raggedright
{\it One shape parameter:} \\
\texttt{ph <a>} {\it proportional hazard} \\
\texttt{wang <lam>} {\it Wang / normal shift} \\
\texttt{dual <b>} {\it dual moment} \\
\texttt{tvar <p>} {\it tail VaR at} $p$ \\
\texttt{ccoc <r>} {\it constant cost of capital} \\
\texttt{cll <b>} {\it capped log-linear} \\
\texttt{clin <slope>} {\it capped linear} \\
\texttt{lep <r>} {\it leverage-equivalent pricing} \\
\texttt{ly <r>} {\it linear yield}

\medskip{\it Several parameters:} \\
\texttt{bitvar <p0> <p1> <w1>} {\it weight} $w_1$ {\it on} $p_1$ \\
\texttt{beta <a> <b>} {\it \rm a\it{} in} $(0, 1]$ \\
\texttt{power <x0> <x1> <alpha>}
\end{minipage}};
\node[fancytitle] at (box.north west) {2. Kinds \& parameters};
\end{tikzpicture}

\columnbreak

%--- Combinators ---
\begin{tikzpicture}
\node [mybox] (box){\begin{minipage}{0.3\textwidth}\raggedright
\texttt{distortion <NAME> minimum dist.<A> dist.<B> ...} \\
\texttt{distortion <NAME> mixture dist.<A> dist.<B> ... wts [<W>]}

\medskip\it \texttt{minimum} $=$ pointwise min of the referenced distortions (min of concaves is concave, so the result is a distortion); \texttt{mixture} $=$ weighted average, weights via \texttt{wts}. Both take distortion \rm references\it, not numbers.
\end{minipage}};
\node[fancytitle] at (box.north west) {3. Combinators};
\end{tikzpicture}

%--- No DecL form ---
\begin{tikzpicture}
\node [mybox] (box){\begin{minipage}{0.3\textwidth}\raggedright
\texttt{wtdtvar} {\it (weighted TVaR) has \rm no\it{} DecL number-list form}: it takes parameter \rm vectors\it{} of knots and weights. Build it in Python, or reach it through the combinators.

\medskip\rm\texttt{Distortion.random\_distortion(n\_knots=5)} \\
\texttt{Distortion.mixture([d1, d2], [.6, .4])}

\medskip\it List the registered kinds with \rm\texttt{Distortion.available\_distortions()}.
\end{minipage}};
\node[fancytitle] at (box.north west) {4. Not declarable};
\end{tikzpicture}

\columnbreak

%--- Usage ---
\begin{tikzpicture}
\node [mybox] (box){\begin{minipage}{0.3\textwidth}\raggedright
\it Apply a distortion to price an aggregate or portfolio:

\medskip\rm\texttt{a = build('agg A 100 claims sev lognorm 50 cv 2 poisson')} \\
\texttt{d = build('distortion D ph 0.9')} \\
\texttt{a.price(0.99, d)} \\
\texttt{p.price(0.99, d)} {\it total \rm and\it{} allocation by unit}

\medskip\it Calibrate to a target instead of guessing a shape: \rm\texttt{p.calibrate\_distortions(coc=0.1, p=0.99)}.

\medskip\it A distortion satisfies $g(0)=0$, $g(1)=1$, increasing and concave; the dual is $g^*(x) = 1 - g(1-x)$; the distorted price is $\int_0^\infty g(S(x))\,dx$.
\end{minipage}};
\node[fancytitle] at (box.north west) {5. Usage};
\end{tikzpicture}

%--- Bounds ---
\begin{tikzpicture}
\node [mybox] (box){\begin{minipage}{0.3\textwidth}\raggedright
\it The reverse question --- \rm which\it{} distortions are consistent with an observed price --- is not a declaration but an analysis. See the \texttt{*Bounds} card: its answers come back as \texttt{bitvar} distortions, which \rm are\it{} declarable.

\medskip\rm\texttt{ab = p.allocation\_bounds(p=0.99)} \\
\texttt{ab.bitvars(1200)} \\
\texttt{dist D bitvar 0.5 0.99 0.3}
\end{minipage}};
\node[fancytitle] at (box.north west) {6. And back again};
\end{tikzpicture}

\makefooter
\end{multicols*}

\end{document}
